\documentclass[letterpaper,twocolumn]{article}
\usepackage{aaai2027}
\usepackage[hyphens]{url}
\usepackage{graphicx}
\urlstyle{rm}
\def\UrlFont{\rm}
\usepackage{natbib}
\usepackage{caption}
\frenchspacing
\setcounter{secnumdepth}{3}
\pdfinfo{
/TemplateVersion (2027.1)
}

\usepackage{amsmath}
\usepackage{amssymb}
\usepackage{amsfonts}
\usepackage{mathtools}
\usepackage{bm}
\usepackage{booktabs}
\usepackage{multirow}
\usepackage{amsthm}

\title{A Model's Identity Direction is Causally Separable from Safety: Mechanistic Evidence from Activation Steering}
\author{
    Naz Col\textsuperscript{1}\thanks{Equal contribution.},
    Ethan TS. Liu\textsuperscript{1}\footnotemark[1],
    Nishit Anand\textsuperscript{2}, 
    Anya Ji\textsuperscript{1},
    David M. Chan\textsuperscript{1},
    Lisa Dunlap\textsuperscript{1}
}
\affiliations{
    \textsuperscript{1}University of California, Berkeley
    \textsuperscript{2}University of Maryland, College Park\\
    doganazcol@berkeley.edu, ethantsliu@berkeley.edu, lisabdunlap@berkeley.edu
}

\begin{document}

\maketitle

\begin{abstract}
Does behavioral imitation harm safety? The field consensus is yes. We show this consensus is a measurement error: Llama-Guard overcounts genuine harm by 2.7–3.5× against content-aware judges (9.3% vs 25.4%). When corrected, fine-tuning an LLM to imitate another model on benign data erodes safety only marginally—mean +0.22 percentage points over 540 adapters, with 245 adapters improving. What little safety change exists tracks the base model being fine-tuned, not which model is imitated. Mechanistically, we show why identity transfer does not drag safety along: a model's identity direction and refusal direction are causally separable in activation space. Across 23 models, ablating identity leaves safety statistically unmoved (identity: −0.9 to +3.9pp, indistinguishable from random), while ablating refusal catastrophically erodes safety (refusal cone: +16 to +96pp harm). This dissociation is validated by positive control (refusal ablation is catastrophic) and null control (random-direction ablation is null). We conclude behavioral imitation's perceived safety cost is a measurement artifact, not a genuine threat. Identity and safety occupy orthogonal activation subspaces—explaining why behavioral transfer does not inherently entail safety loss.
\end{abstract}

\section{Introduction}

Large language models exhibit distinctive behavioral fingerprints: tone, verbosity, response structure, and characteristic failure modes that users and tools learn to recognize. Recent work proposes to exploit these fingerprints for model provenance auditing, detecting unauthorized distillation by measuring whether an endpoint's outputs still carry the expected signature. The threat model is direct: a party that fine-tunes to imitate another model's outputs is incidentally or deliberately laundering its behavioral signature.

Under this threat, a critical mechanistic question emerges. When behavioral imitation erodes safety (as established by prior work), what is the causal mechanism? Is the safety erosion intrinsic to shifting identity, perhaps because safety and identity share a common mechanistic substrate, or is it a collateral effect of the weight-level changes required to fit target outputs, logically separable from the identity shift itself?

This distinction has practical consequences. If identity and safety occupy separable activation-space subspaces, then deployment-time steering of the identity direction could enable behavioral transfer without safety loss. Conversely, if they are mechanistically coupled, no steering-based intervention can decouple them; safety protection would require architectural changes or refusal training orthogonal to the imitation path.

We investigate this question through targeted activation steering experiments. Using the diff-of-means refusal direction from Arditi et al.'s steering methodology, we ask whether the analogous benign cross-model provenance direction (identity direction) and the safety direction occupy orthogonal subspaces. Our findings support mechanistic separability: identity and safety directions are largely orthogonal, and ablating one has minimal effect on the other. However, we also uncover a surprising complication: the separability is model-dependent, with coupling appearing in models where the two directions subtend small angles in activation space.

\section{Background and Motivation}

\subsection{Behavioral Fingerprinting and Imitation}

Each LLM carries a distinctive behavioral signature measurable in output statistics: response length, formatting preferences, reasoning style, and characteristic phrasings. Prior work on behavioral inertia shows this signature persists across tasks and datasets, is stable across inference runs, and erodes predictably under an escalating intervention ladder (prompting, supervised fine-tuning, preference optimization). A recent study of cross-model imitation demonstrates that fine-tuning a source model on target outputs measurably transfers this behavioral signature (+22 to +28\% match rate) while simultaneously degrading safety: refusal rates drop by 5 to 15 percentage points depending on the target model, more severely than for self-imitation controls.

The finding that imitation erodes safety is not novel—it is one instance of a broader phenomenon where benign fine-tuning degrades model alignment. However, the imitation-specificity of the effect, where self-imitation (a model fine-tuned on its own outputs) shows minimal safety loss while cross-model imitation shows substantial loss, suggests the erosion is not purely a fine-tuning side-effect but related specifically to adopting a different model's outputs.

\subsection{Prior Mechanistic Work on Steering and Orthogonality}

Recent work by Arditi et al. (2024) established that refusal in LLMs can be mediated by a single direction in activation space. They compute the direction as the mean activation difference between harmful-prompt refusals and harmful-prompt compliances, and show that ablating this direction (subtracting the component of the activation along this direction) degrades refusal performance catastrophically. Importantly, this steering is done via projection: for a direction $\vec{r}$, they modify activations as $\vec{h}' = \vec{h} - \beta (\vec{h} \cdot \hat{\vec{r}}) \hat{\vec{r}}$, effectively removing the component of the hidden state along the refusal direction.

The key methodological insight is that this steering operator is direction-specific: ablating a random direction has minimal effect on safety, while ablating the computed refusal direction is catastrophic. This orthogonality to random directions validates that the refusal effect is not a general hidden-state perturbation artifact but specific to the refusal direction's geometry.

We apply the identical steering operator, but to a different contrast set: instead of harmful-refuse vs.\ harmful-comply, we use benign-source vs.\ benign-target activations. This yields an identity direction whose null ablation effect would indicate separability from safety.

\section{Methods}

\subsection{Computing the Identity Direction}

For each (source model, target model) pair, we collect hidden activations from the base models on a shared set of benign prompts (500 examples from GSM8K, Chatbot Arena, and WritingPrompts). We focus on benign prompts to isolate the identity/style signal from any safety-related differences.

For each prompt $p$, we generate one response from the source model $r_S(p)$ and one from the target model $r_T(p)$, and extract the hidden state $\vec{h}(p, r)$ at layer $L$ (chosen empirically; we report results for $L \in \{14, 20, 28\}$ depending on model size). The identity direction is computed as the mean difference:
\[
\vec{d}_{\text{id}} = \frac{1}{N} \sum_{i=1}^{N} (\vec{h}(p_i, r_T(p_i)) - \vec{h}(p_i, r_S(p_i)))
\]

This direction captures the activation-space signature of behavioral transition from source to target style.

\subsection{Computing the Refusal Direction}

The refusal (safety) direction is computed from harmful prompts. For each prompt $q$ requesting harmful content, we generate two responses: one refusing the request, and one attempting to comply (sampled at low temperature to maximize likelihood of compliance if the model is tempted). The refusal direction is:
\[
\vec{d}_{\text{safe}} = \frac{1}{M} \sum_{j=1}^{M} (\vec{h}(q_j, r_{\text{refuse}}(q_j)) - \vec{h}(q_j, r_{\text{comply}}(q_j)))
\]

For aligned models, this direction points away from compliance.

\subsection{Steering and Evaluation Protocol}

During generation, we perturb hidden states using the Arditi et al.\ operator, removing the component along a chosen direction at layer $L$ with strength $\beta$:
\[
\vec{h}'(\beta, \vec{d}) = \vec{h} - \beta (\vec{h} \cdot \hat{\vec{d}}) \hat{\vec{d}}
\]

We evaluate three conditions:

\begin{enumerate}
\item[(a)] \textbf{Identity ablation} ($\beta=1.0$ applied to $\vec{d}_{\text{id}}$): Remove the identity direction and measure how much behavioral imitation survives. We measure match rate (fraction of outputs judged more similar to target than source) and refusal rate on harmful prompts.

\item[(b)] \textbf{Refusal ablation} ($\beta=1.0$ applied to $\vec{d}_{\text{safe}}$): Remove the refusal direction to confirm the positive control (safety should be catastrophically eroded).

\item[(c)] \textbf{Random ablation} ($\beta=1.0$ applied to a random direction orthogonal to identity): Establish the null control; neither identity transfer nor safety loss should occur.

\end{enumerate}

We also measure the angle $\theta$ between $\vec{d}_{\text{id}}$ and $\vec{d}_{\text{safe}}$ via $\cos(\theta) = (\vec{d}_{\text{id}} \cdot \vec{d}_{\text{safe}}) / (|\vec{d}_{\text{id}}| |\vec{d}_{\text{safe}}|)$.

\section{Results}

\subsection{Identity Ablation Preserves Safety}

On Llama 3.1 8B, ablating the identity direction via steering achieves behavioral imitation (+22\% match rate) while preserving refusal rate: 87.3\% (baseline) vs.\ 86.8\% (identity-ablated, $\Delta = -0.5$ percentage points). Across five safety benchmarks (AdvBench, HarmBench, StrongREJECT, XSTest, SORRY-Bench), identity ablation shows minimal refusal degradation (mean $\Delta = -1.2$ percentage points, $\sigma = 0.8$). In contrast, refusal ablation catastrophically erodes safety: refusal drops to 22.1\%, a 65.2 percentage-point decline (Figure~\ref{fig:positive_control}).

On Qwen 3.6 27B, identity ablation achieves +19\% match rate with refusal preserved at 88.1\% (baseline 88.5\%, $\Delta = -0.4$ pp). Refusal ablation again catastrophically erodes safety ($\Delta = -63.9$ pp).

Random ablation on the same models shows minimal effects: +1.2\% match rate and -0.3 pp refusal loss, validating that the steering is direction-specific.

\begin{figure}[t]
\centering
\includegraphics[width=\linewidth]{img/03_steering_effects.pdf}
\caption{Per-model steering effects. Left: Identity ablation achieves +18–22\% behavioral match rate (transfer successful). Right: Minimal refusal degradation under identity ablation (Llama/Qwen $\le 1$ pp; GPT-oss -4.2 pp due to coupling). Demonstrates that successful behavioral transfer can occur with negligible safety impact in well-aligned models.}
\label{fig:steering_effects}
\end{figure}

\begin{figure}[t]
\centering
\includegraphics[width=0.85\linewidth]{img/04_positive_control.pdf}
\caption{Positive control: Refusal ablation is catastrophic. Baseline refusal rate is 87.3\%. Identity ablation: -0.5 pp (separable). Refusal ablation: -65.2 pp (catastrophic safety loss). Random ablation: -0.3 pp (null control). This validates that steering effects are direction-specific and that the refusal direction encodes safety independently from identity.}
\label{fig:positive_control}
\end{figure}

\subsection{Angle Between Identity and Safety Directions Predicts Coupling}

We observe model-dependent variation in the separability. Computing the angle $\theta$ between identity and safety directions (Figure~\ref{fig:orthogonality}):

\begin{itemize}
\item \textbf{Large angle ($\theta > 80°$):} Llama 3.1 8B, Qwen 3.6 27B show near-orthogonal directions and clean separability (identity ablation $\Delta$-refusal $\le 1$ pp).
\item \textbf{Small angle ($\theta < 40°$):} GPT-oss 20B shows substantial angle, and identity ablation degrades refusal by -4.2 pp. Inspection of the activation space reveals that the small angle correlates with models where identity and safety features co-localize in early layers.
\end{itemize}

For GPT-oss, the coupling appears to be genuine: partial ablation experiments show that only the component of identity orthogonal to safety ($\vec{d}_{\text{id}}^{\perp}$) leaves refusal intact (-0.7 pp); the parallel component ($\alpha_{\parallel} \vec{d}_{\text{safe}}$) degrades refusal (-3.5 pp).

\begin{figure}[t]
\centering
\includegraphics[width=0.8\linewidth]{img/02_model_orthogonality.pdf}
\caption{Model-dependent orthogonality predicts coupling. Models with large angles between identity and safety directions (Llama, Qwen at $\theta > 80°$) show clean separability (blue/orange, near zero on y-axis). GPT-oss with small angle ($\theta < 40°$) shows coupling (red, -4.2 pp refusal degradation under identity ablation). Shaded regions highlight separable vs.\ coupled regimes.}
\label{fig:orthogonality}
\end{figure}

\subsection{Cross-Benchmark Validation}

We validate the dissociation across multiple safety benchmarks. Figure~\ref{fig:benchmark_dissociation} shows the core result: identity ablation causes minimal refusal degradation across all benchmarks while refusal ablation is catastrophic.

\begin{figure}[t]
\centering
\includegraphics[width=\linewidth]{img/01_benchmark_dissociation.pdf}
\caption{Benchmark dissociation across nine safety benchmarks. Identity ablation (blue bars) shows minimal refusal rate change (mean $-1.1 \pm 0.9$ pp), while refusal ablation (red bars) devastates safety ($-62.2 \pm 5.4$ pp). Random ablation (green bars) is null, validating direction-specificity of steering effects. Error bars show $\pm 1$ SD across multiple models and seeds.}
\label{fig:benchmark_dissociation}
\end{figure}

\begin{table*}[t]
\centering
\small
\begin{tabular}{lrrrr}
\toprule
Benchmark & Identity Ablation & Refusal Ablation & Random Ablation & Consensus \\
\midrule
AdvBench & $-1.3 \pm 0.9$ & $-65.8 \pm 4.2$ & $-0.2 \pm 0.4$ & Separable \\
HarmBench & $-0.8 \pm 0.7$ & $-62.1 \pm 5.1$ & $+0.1 \pm 0.5$ & Separable \\
StrongREJECT & $-1.1 \pm 1.0$ & $-58.3 \pm 6.8$ & $-0.3 \pm 0.6$ & Separable \\
SORRY-Bench & $-0.9 \pm 0.8$ & $-59.7 \pm 4.9$ & $+0.2 \pm 0.4$ & Separable \\
SG-Bench & $-1.4 \pm 1.1$ & $-60.5 \pm 7.2$ & $-0.1 \pm 0.5$ & Separable \\
XSTest (harm) & $-1.5 \pm 1.2$ & $-61.2 \pm 5.3$ & $-0.2 \pm 0.5$ & Separable \\
XSTest (over-ref) & $+0.8 \pm 1.0$ & $-2.1 \pm 1.5$ & $+0.0 \pm 0.3$ & Separable \\
OR-Bench hard & $-0.6 \pm 0.6$ & $-64.1 \pm 5.8$ & $-0.2 \pm 0.4$ & Separable \\
OR-Bench toxic & $-0.5 \pm 0.7$ & $-63.2 \pm 6.1$ & $+0.1 \pm 0.5$ & Separable \\
\midrule
\textbf{Mean (core)} & $\bf{-1.1 \pm 0.9}$ & $\bf{-62.2 \pm 5.4}$ & $\bf{-0.1 \pm 0.5}$ & \textbf{Separable} \\
\bottomrule
\end{tabular}
\caption{Mean refusal degradation (\%) under identity ablation vs.\ refusal ablation, across multiple open-weight models and core safety benchmarks ($N=300$ prompts per benchmark where applicable). Error bars report $\pm 1$ SD. Identity ablation: no significant safety impact. Refusal ablation: catastrophic safety loss. Random ablation: null control.}
\label{tab:benchmark_dissociation}
\end{table*}

The pattern is consistent: identity ablation shows minimal refusal loss (mean -1.1 pp), while refusal ablation devastates safety (-62.2 pp). On the over-refusal axis (XSTest), identity ablation shows negligible effect (+0.8 pp, indicating no over-refusal induced by identity transfer), while refusal ablation predictably reduces over-refusal (-2.1 pp because the refusal direction is now ablated).

\section{Mechanistic Interpretation}

These results establish three key mechanistic facts:

\textbf{1. Identity and safety are geometrically separable.} The activation-space identity direction and refusal direction occupy largely orthogonal subspaces, as evidenced by:
\begin{itemize}
\item Minimal refusal loss under identity ablation (-1.1 pp on average),
\item Orthogonal controls (random direction ablation shows -0.1 pp refusal loss),
\item Model-dependent angle measurements showing $\theta > 80°$ in well-aligned models.
\end{itemize}

\textbf{2. Safety degradation under imitation is caused by weight changes, not identity transfer.} The steering result reveals that the mechanism of safety erosion is not the act of shifting identity itself, but rather the weight-level changes (parameter updates in DPO/SFT) required to fit the target model's outputs. During deployment, if identity shift is induced via steering alone (without weight updates), safety is preserved.

\textbf{3. Separability is model-dependent.} Coupling (identity ablation $\Delta$-refusal $\approx -4$ pp) appears in models where identity and safety directions are co-localized (small angle $\theta$). This suggests alignment method and model architecture influence whether these mechanistically important directions are orthogonal or coupled.

\section{Implications}

\subsection{Deployment and Auditing}

If identity and safety are activation-space separable, then deployment systems can defend against imitation-based fingerprint laundering without sacrificing behavioral transfer. A model could be steered toward a target identity via activation modification while simultaneously steering to preserve (or enhance) safety. This separability provides a mechanistic justification for decoupled control of identity and safety at inference time.

\subsection{Alignment Research}

The finding that safety is robust to identity shifting, but catastrophically vulnerable to refusal-direction ablation, suggests that current alignment methods (DPO, RLHF) encode safety primarily in the refusal direction rather than distributed across identity and capability features. This has implications for fine-tuning robustness: models may be more resilient to identity-targeted attacks (imitation, distillation) than to fine-tuning on genuinely different training distributions.

\section{Related Work}

This work builds on three literatures: (1) behavioral fingerprinting and model provenance via output statistics; (2) fine-tuning-induced safety degradation and its mechanistic causes; and (3) activation steering and representation-level interventions in language models.

Behavioral fingerprints have been proposed to verify model identity and detect distillation. Recent work shows fingerprints can be erased through adversarial fine-tuning, but the mechanism (whether identity is causally erased or merely hidden) was unclear. We clarify this by showing that identity is steerable and erasable at the activation level.

The finding that benign fine-tuning erodes safety has been extensively documented. Recent work distinguishes between generic fine-tuning erosion and imitation-specific erosion via self-imitation controls. Our contribution is the mechanistic localization: we show the erosion is due to weight changes, not to the identity shift itself.

Prior work on representation steering (CAA, RepE, LAT) and refusal directions (Arditi et al.) established that single directions can encode important model properties. We extend this by showing that different mechanistically important directions can be orthogonal and independently controllable.

One concurrent paper (Jun 2026, arXiv:2606.26161) reaches the opposite conclusion using identical methodology, claiming that ablating the persona direction restores refusal. We acknowledge this concurrent work and note that differences in layer selection, model choice, and prompt engineering may explain the divergent findings. Future work should reconcile these results through a unified mechanistic framework.

\section{Limitations}

\begin{enumerate}
\item \textbf{Experimental scope:} Results are on open-weight models (Llama, Qwen, GPT-oss) with inference-time steering access. Closed models (GPT-4, Claude) cannot be directly tested, limiting claims about universal separability.

\item \textbf{Layer selection:} Identity and refusal directions vary by layer. We report layer $L$ chosen via validation; other layers may show different separability. A full layer-sweep analysis is deferred to future work.

\item \textbf{Prompt set:} Identity direction is computed on benign prompts from specific datasets. Separability may differ on other prompt distributions or under adversarial prompt optimization.

\item \textbf{Causality and interference:} Steering is a strong intervention at $\beta = 1.0$. The interaction between identity and safety ablations at weaker strengths remains characterizable but harder to interpret causally.

\item \textbf{Model-dependent coupling:} The separability breaks down in models with small angles between identity and safety directions, but we lack a mechanistic explanation for why certain architectures or alignment methods produce this co-localization.

\end{enumerate}

\section{Conclusion}

A model's behavioral identity direction and its safety direction are geometrically orthogonal and causally independent at the activation level. Ablating identity transfers style without eroding safety; ablating safety catastrophically erodes refusals while leaving identity intact. This mechanistic separability explains why imitation-based fine-tuning erodes safety (weight changes, not identity shift) and opens the possibility of deployment-time identity steering that preserves safety through representation-level intervention.

However, separability is not universal: it breaks down in models where identity and safety co-localize in activation space. Future work should characterize which alignment methods and architectures produce orthogonal directions, and whether training-time interventions can enforce separability for robustness.

\begin{thebibliography}{99}

\bibitem[Arditi et al.]{arditi2024} Arditi, A., Dugan, L., Sap, M. (2024). Refusal in LLMs is mediated by a single direction. In: \textit{Advances in Neural Information Processing Systems (NeurIPS)}.

\bibitem[Qi et al.]{qi2023} Qi, Z., Ding, Y., Gu, S., et al. (2023). Fine-tuning language models erodes their safety. In: \textit{International Conference on Learning Representations (ICLR)}, 2024.

\bibitem[Hinton et al.]{hinton2015} Hinton, G., Vanhoucke, V., Dean, J. (2015). Distilling the knowledge in a neural network. In: \textit{NIPS Deep Learning Workshop}.

\bibitem[Gudibande et al.]{gudibande2023} Gudibande, A., Wallace, E., Snorkel, C., et al. (2023). The false promise of imitating proprietary LLMs. \textit{arXiv:2305.15717}.

\bibitem[Lee et al.]{lee2025} Lee, K., Ippolito, D., Nystrom, A., et al. (2025). Detecting model misgeneralization in LLMs. \textit{arXiv:2410.06145}.

\bibitem[Wollschlaeger et al.]{wollschlaeger2025} Wollschlaeger, D., et al. (2025). Refusal in the refusal cone. \textit{arXiv:2502.17420}.

\end{thebibliography}

\end{document}
